\documentclass[11pt, oneside]{amsart}

\usepackage{geometry}
\geometry{letterpaper}
\usepackage{ned-common}
\usepackage{ned-calculus}
\usepackage{ned-linear-algebra}

\begin{document}
\title{Root Mean Square}
\maketitle

\section{Root Mean Square Voltage}

\begin{enumerate}
  \item Say you attach a purely resistive circuit to an AC power source.

  \item We will be assuming that AC is a sinusoidal running at angular
  frequency $\omega$. That is: $v(t) = V_\text{max} \cos \omega t$.

  \item The \define{instantaneous power} dissipated is $p(t) = v(t) i(t)
  = v^2(t) / R = V^2_\text{max} \cos^2{\omega t} / R$.

  \item We can also define the \define{average power}. This is

  \begin{nedqn}
    P
  \eqcol
    \frac{1}{T}
    \int_0^T v^2(t) / R \dt
  \\
  \eqcol
    \frac{1}{T}
    \int_0^T
      V^2_\text{max} \cos^2{\omega t}
      / R
      \dt
  \end{nedqn}

  \noindent
  This is also often called the \define{real power}, and also
  \define{active power}.

  \item Can we simplify the expression for real power? We can't replace
  $v^2(t)$ with $V_\text{AVG}^2$, since $V_\text{AVG} = 0$!

  \item We also cannot say $P = V_\text{max}^2 / R$; that's the peak
  \emph{instantaneous} power, but almost everywhere less power is
  consumed!

  \item Let's imagine a non-sinusoidal voltage source. Consider a 1V DC
  source attached for 2sec to a $1\Omega$ circuit. This gives
  instantaneous power of $p(t) = 1 \text{W}$. Next, consider a 2V DC
  source attached for 1sec, and then 0V for 1sec. This gives $p(t) = 4
  \text{W}$ for 1sec, and then $p(t) = 0\text{W}$ for 1sec. Even though
  both examples have the same average voltage applied, the second
  circuit draws 2W of real (average) power over the 2sec cycle, versus
  1W.

  \item However, we may define the \define{root mean square} voltage:

  \begin{nedqn}
    V_\text{RMS}
  \eqcol
    \sqrt{\frac{1}{T} \int_0^T v^2(t) \dt}
  \intertext{That will allow us to write (for resistive circuits, at least):}
    P
  \eqcol
    V_\text{RMS}^2 / R
  \end{nedqn}

  \item The RMS voltage is the voltage that, if applied continuously to
  the (purely resistive) circuit, would draw as much real power as the
  alternating voltage supply.

  \item Per our earlier example of 1sec 2V and 1sec 0V: $V_\text{RMS} =
  \sqrt{2}\text{V}$.

  \item Let's return to our typical case where the circuit is driven by
  a sinusoidal voltage running at angular frequency $\omega$. Then I
  argue

  \begin{nedqn}
    V^2_\text{RMS}
  \eqcol
    \frac{1}{T}
    \int_0^T V^2_\text{max} \cos^2\omega t \dt
  \\\eqcol
    0.5 V^2_\text{max}
  \end{nedqn}

  \noindent
  This follows from the typical identity about symmetry of
  sinusoids/circles.
\end{enumerate}

\section{Resistance}

\begin{enumerate}
  \item We already know about resistance, denoted $R$ and measured in
  Ohms ($\Omega$). Through a pure resistor, $V = IR$. The voltage drop
  across a pure resistor is proportional to the magnitude of the
  current.
\end{enumerate}

\section{Inductors}

\begin{enumerate}

  \item An \define{inductor} ``opposes'' change in current.

  \item When current (magnitude) is rising, energy is stored in the
  magnetic field. Electromagnetic induction is the property that an
  increasing current magnitude will store energy in the magnetic field
  (and a decreasing current magnitude will consume stored energy).

  \item Because the energy stored in the magnetic field must be paid
  for, it is natural that there is a voltage drop across the inductor
  when the inductor is charging. It gives a voltage rise across the
  inductor when the inductor is discharging. This voltage exists because
  the change to the magnetic field always happens hand-in-hand with a
  corresponding change in the \emph{electric field}.

  \item When current is zero, the inductor has stored no energy in the
  magnetic field. We don't say that the magnetic flux through the
  inductor core is zero (maybe there some external EMF), only that the
  inductor hasn't stored any energy there. The inductor charges no
  matter which direction you start flowing current across it. It will
  discharge this energy as you slow current and return it to zero
  current.

  \item You \emph{might} be tempted to think of the inductor charging as
  posing a \emph{resistance} to current flow through the inductor. But
  don't do that! Resistance taxes more energy per charge the faster you
  push charges through (proportional to current). For an inductor, the
  energy taxed for pushing charge through is constant for any level of
  current, for a given rate of current change. Literally, you are
  pushing charge through a voltage (electric field) that opposes its
  motion.

  \item When current stops rising, energy stops being stored in magnetic
  field, and voltage drop across inductor is zero. As long as the
  current remains constant, no energy stored in the field will be
  ``consumed.'' The energy stored remains constant, ready to be returned
  when current begins to slow.

  \item When you begin to decrease the current, the inductor begins to
  deliver the power stored in the magnetic field, which is a voltage
  \emph{increase} (rather than drop) across the inductor. Again, do not
  think of this as a ``negative resistance,'' since the energy ``push''
  given each unit charge is constant, and not proportional to current
  magnitude. It really is a voltage that is induced.

  \item Note that energy is stored no matter whether you rise from zero
  current or fall to ``negative'' current. This is just a matter of
  orientation. It's really the current \emph{magnitude} that matters. Of
  course, the charging and discharging push direction swaps orientation.
  But you can't fall below ``zero energy stored.''

  \item At the same current magnitude, always the same amount of energy
  has been stored in the field. You can raise current slowly over a long
  time, which increases stored energy slowly. But then if you try to
  bring current back down quickly, the stored energy will ``resist''
  harder, but for a shorter period of time.

  \item The Henry (H) is equivalent to Volts (of back-EMF) per ampere
  per second (rate of change). The symbol $L$ is for Heinrich Lenz, but
  the unit is for Joseph Henry.

  \item Anyway, the voltage drop across an inductor is:

  \begin{nedqn}
    v(t)
  \eqcol
    L \fpartial[i(t)]{t}
  \end{nedqn}

  \item The energy stored in an inductor is:

  \begin{nedqn}
    U
  \eqcol
    \int_0^T v(t) i(t) \dt
  \\\eqcol
    \int_0^T L \fpartial[i(t)]{t} i(t) \dt
  \\\eqcol
    \int_0^I L i \diff{i}
  \nedcomment{change of variables}%
  \\\eqcol
    \frac{1}{2} L I^2
  \end{nedqn}

  \noindent
  Maybe be careful to note that change-of-variables only works if
  current started at zero at time zero.
\end{enumerate}

\section{Electromagnetic Induction}

\begin{enumerate}
  \item I don't want to get \emph{too} deep into induction. You don't
  need to understand how resistors resist in order to use a resistor.
  You don't need to know how to make an inductor to use one.

  \item Still, this is an important concept. The fundamental idea is
  the Maxwell-Faraday equation:

  \begin{nedqn}
    \grad \times \mtxE
  \eqcol
    - \fpartial[\mtxB]{t}
  \end{nedqn}

  \item $\grad \times \mtxE$ is the ``curl'' of $\mtxE$, sometimes
  written $\operatorname{rot} \mtxE$ or $\operatorname{curl} \mtxE$. It
  maps each point in the electric field $\mtxE$ to a vector about which
  the field rotates at that point. The length of this vector is
  proportional to the rate of the rotation.

  \item $-\fpartial[\mtxB]{t}$ is the rate at which the magnetic field
  is changing. This is a map from vectors to vectors, since the magnetic
  field at each point has both a strength and an orientation.

  \item Basically, the Maxwell-Faraday equation is saying that in order
  for the electric field to change, the magnetic field must change. And
  vice versa.

  \item Even in vacuum, changes to magnetic and electric fields must
  correspond. Magnetic and electric fields, even in vacuum, store energy
  simply in the field itself; increasing the strength of a magnetic
  (equivalently, an electric) field ``takes'' energy, which is stored in
  the fields. Weakening the field will ``take'' energy out of the
  fields, which must flow elsewhere.

  \item Changing
\end{enumerate}

\end{document}
